---
format: yohaku/0.2
target: text
macro: '[[macros/raw]]'
---

## 正文
````text
"""Offline frozen-document retriever. Python 3 standard library only."""
import argparse
from bisect import bisect_right
import hashlib
import types
import json
import os
from pathlib import Path
import sys
import tempfile

ROOT = Path(__file__).resolve().parent
BASELINE_SEAL = "5ccd119bd634f72630239a1b9787e3aab4da28622a6ee397b527dd97a9341b7e"
SNAPSHOT_SHA = "893a261990e2c4d9b5fec9043fee872b9bf510e4f8db8451ddfc88c1a21686ae"
FIELDS = ("path", "title", "body")
STATUSES = ("exists", "missing", "unknown")

def sha(raw):
    return hashlib.sha256(raw).hexdigest()

def load_dataset(root=ROOT):
    """Check trusted fixed bytes before importing the unchanged local verifier."""
    base = Path(root) / "baseline"
    raw = (base / "BASELINE.json").read_bytes()
    if sha(raw) != BASELINE_SEAL:
        raise ValueError("baseline seal hash mismatch")
    seal = json.loads(raw)
    verified_bytes = {}
    for name, row in seal["files"].items():
        raw = (base / name).read_bytes()
        if sha(raw) != row["sha256"] or len(raw) != row["bytes"]:
            raise ValueError("baseline asset hash mismatch: " + name)
        verified_bytes[name] = raw
    # Execute only the hashed source bytes; never read or write a .pyc cache.
    verifier = types.ModuleType("_frozen_material_verifier")
    exec(compile(verified_bytes["verify_materials.py"], str(base / "verify_materials.py"), "exec"), verifier.__dict__)
    verifier.verify_bundle(base / "materials.zip", external_dir=base, input_sha256=SNAPSHOT_SHA)
    data = json.loads(verified_bytes["snapshot.json"])
    index = json.loads(verified_bytes["index.json"])
    links = json.loads(verified_bytes["links.json"])
    rows = {d["source"]: d for d in index["documents"]}
    documents = []
    for doc in data["documents"]:
        row = rows[doc["source"]]
        documents.append(dict(source=doc["source"], source_version=doc["source_version"],
                              title=row["title"], body=doc["body"]))
    return documents, links["links"]

def link_counts(documents, links):
    counts = {d["source"]: dict.fromkeys(STATUSES, 0) for d in documents}
    for link in links:
        if link["source"] not in counts or link["status"] not in STATUSES:
            raise ValueError("invalid link source or status")
        counts[link["source"]][link["status"]] += 1
    return counts

def snippet(body, needle=None):
    """Return exact original complete lines, including original line endings."""
    lines = body.splitlines(keepends=True)
    if not lines:
        return "", 0, 0
    if needle is None:
        first = next((i for i, line in enumerate(lines) if line.strip()), None)
        if first is None:
            return "", 0, 0
        stop = first + 1
        while stop < len(lines) and lines[stop].strip():
            stop += 1
        return "".join(lines[first:stop]), first + 1, stop
    folded = body.casefold()
    hit = folded.find(needle)
    if hit < 0 or not needle:
        raise ValueError("snippet requested without body match")
    # Casefold may expand a character (e.g. ß -> ss). Translate both ends,
    # not merely the start, back to original Unicode code-point offsets.
    offsets = [0]
    for char in body:
        offsets.append(offsets[-1] + len(char.casefold()))
    begin = bisect_right(offsets, hit) - 1
    last = bisect_right(offsets, hit + len(needle) - 1) - 1
    ends = []
    total = 0
    for line in lines:
        total += len(line)
        ends.append(total)
    first_line = bisect_right(ends, begin)
    last_line = bisect_right(ends, last)
    return "".join(lines[first_line:last_line + 1]), first_line + 1, last_line + 1

def search(documents, links, field, query, limit=None):
    if field not in FIELDS + ("all",):
        raise ValueError("invalid field")
    if not isinstance(query, str) or not query.strip():
        raise ValueError("query must contain a non-whitespace character")
    if limit is not None and (type(limit) is not int or limit <= 0):
        raise ValueError("limit must be a positive integer")
    needle = query.casefold()
    selected = FIELDS if field == "all" else (field,)
    counts = link_counts(documents, links)
    results = []
    for doc in sorted(documents, key=lambda d: d["source"]):
        values = {"path": doc["source"], "title": doc["title"], "body": doc["body"]}
        matched = [key for key in selected if needle in values[key].casefold()]
        if not matched:
            continue
        text, start, end = snippet(doc["body"], needle if "body" in matched else None)
        results.append(dict(source=doc["source"], source_version=doc["source_version"],
                            title=doc["title"], matched_fields=matched, snippet=text,
                            line_start=start, line_end=end, link_counts=counts[doc["source"]]))
    return dict(query=query, field=field, total=len(results),
                results=results if limit is None else results[:limit])

class Parser(argparse.ArgumentParser):
    def error(self, message):
        raise ValueError(message)

def write_output(path, text):
    target = Path(path).resolve()
    base = (ROOT / "baseline").resolve()
    if target == base or base in target.parents:
        raise ValueError("output must not modify read-only baseline")
    # Atomic replace: no partial JSON after a failed write.
    temporary = None
    try:
        with tempfile.NamedTemporaryFile(mode="w", encoding="utf-8", newline="",
                                         dir=target.parent, delete=False) as f:
            temporary = Path(f.name)
            f.write(text)
        os.replace(temporary, target)
    finally:
        if temporary is not None and temporary.exists():
            temporary.unlink()

def main(argv=None):
    p = Parser(description=__doc__, add_help=False, allow_abbrev=False)
    p.add_argument("--help", action="store_true")
    p.add_argument("command", nargs="?", choices=["search"])
    p.add_argument("--field", choices=FIELDS + ("all",))
    p.add_argument("--query")
    p.add_argument("--limit", type=int)
    p.add_argument("--output")
    try:
        args = p.parse_args(argv)
        if args.help:
            print(json.dumps({"usage": "python retriever.py search --field path|title|body|all --query TEXT [--limit N] [--output PATH]"}))
            return 0
        if args.command != "search" or args.field is None or args.query is None:
            raise ValueError("search requires --field and --query")
        if not args.query.strip():
            raise ValueError("query must contain a non-whitespace character")
        if args.limit is not None and args.limit <= 0:
            raise ValueError("limit must be a positive integer")
        documents, links = load_dataset()
        result = search(documents, links, args.field, args.query, args.limit)
        text = json.dumps(result, ensure_ascii=False, indent=2) + "\n"
        if args.output is not None:
            write_output(args.output, text)
        sys.stdout.write(text)
        return 0
    except Exception as exc:
        message = str(exc) or type(exc).__name__
        print(json.dumps({"error": message}, ensure_ascii=False))
        print("ERROR: " + message, file=sys.stderr)
        return 2

if __name__ == "__main__":
    sys.exit(main())
````

## 输出与关联
```yaml
output: "research/workflow-delivery/inherited/retriever_core.py"
```
